\subsection{Radial volumes on scalar Hopf manifolds}
\label{subsec:radial-gauduchon-lck-volumes}

Let $n\geq2$, put $m=n-1$, and consider the scalar Hopf manifold
\[
 H_n=(\mathbb C^n\setminus\{0\})/\langle z\mapsto\mathrm e z\rangle,
 \qquad q:H_n\longrightarrow\mathbb P^m.
\]
Normalize the Fubini--Study form by $\int_{\mathbb P^m}\kappa^m=1$. The circle-valued function $r=\log|z|^2\pmod2$ has a global differential. Set
\[
 \alpha=q^*\kappa=\ddc r,\qquad
 \beta=\mathrm dr\wedge\dc r.
\]
Both forms descend to $H_n$. The form $\alpha$ has rank $m$ and kernel the complex radial line, while $\beta$ has rank one and the complementary horizontal kernel. Thus $\alpha^m\wedge\beta$ is a positive volume form. On each elliptic fibre, $\dc r=\mathrm d\arg z/(2\pi)$, so
\[
 \int_{q^{-1}(x)}\beta=2,\qquad
 \int_{H_n}\alpha^m\wedge\beta=2.
\]
We consider the entire radial family
\begin{equation}\label{eq:radial-lck-family}
 \omega=C(r)\alpha+D(r)\beta,\qquad
 C,D\in C^\infty(\mathbb R/2\mathbb Z),\quad C,D>0.
\end{equation}
The radial circle is oriented by increasing $r$ from zero to two. All Lee periods below use this orientation.

\begin{lemma}\label{lem:radial-lck-gauduchon-equation}
Every metric in \eqref{eq:radial-lck-family} is strictly LCK. Its Lee form and period are
\[
 \lambda=l(r)\,\mathrm dr,\qquad
 l=\frac{C'-D}{C},\qquad
 L=-\int_0^2\frac{D}{C}\,\mathrm dr<0.
\]
It is Gauduchon if and only if $C^{m-1}(C'-D)=k$ for a constant $k$. For a fixed positive smooth two-periodic density $F=C^mD$, this condition is equivalent to
\begin{equation}\label{eq:radial-lck-periodic-ode}
 C^mC'=kC+F,\qquad D=\frac{F}{C^m}.
\end{equation}
Every positive periodic solution has $k<0$, and
\begin{equation}\label{eq:radial-lck-period-identities}
 \begin{gathered}
 k\int_0^2C\,\mathrm dr+\int_0^2F\,\mathrm dr=0,
 \\
 l=kC^{-m},\qquad
 L(C)=-\int_0^2\frac{F}{C^{m+1}}\,\mathrm dr.
 \end{gathered}
\end{equation}
\end{lemma}
\begin{proof}
The differential identities are
\[
 \begin{gathered}
 \alpha^{m+1}=\beta^2=0,\qquad
 \mathrm d\beta=-\mathrm dr\wedge\alpha,\\
 \dc\beta=-\alpha\wedge\dc r,\qquad
 \ddc\beta=-\alpha^2.
 \end{gathered}
\]
They give $\mathrm d\omega=(C'-D)\mathrm dr\wedge\alpha=\lambda\wedge\omega$. This closed form is the actual Lee form, since wedging one-forms with $\omega$ is injective. The integral of $C'/C$ over a period is zero, proving the formula for $L$ and its nonexactness. Two radial Lee forms have the same de Rham class exactly when their periods agree: a zero-period difference has a periodic primitive, and the radial circle detects the converse.

For a smooth radial scalar $G$, the chain rule gives $\ddc G=G'\alpha+G''\beta$. Also $\dc(G\beta)=-G\alpha\wedge\dc r$, so
\[
 \ddc(G\beta)=-G'\alpha\wedge\beta-G\alpha^2.
\]
Expanding the powers of $\omega$ therefore gives
\[
 \begin{aligned}
 \omega^m&=C^m\alpha^m+mC^{m-1}D\alpha^{m-1}\wedge\beta,\\
 \omega^n&=nC^mD\alpha^m\wedge\beta,
 \end{aligned}
\]
and
\[
 \ddc(\omega^m)=
 \bigl[(C^m)''-m(C^{m-1}D)'\bigr]\alpha^m\wedge\beta.
\]
Its vanishing is equivalent to constancy of $C^{m-1}(C'-D)$. Multiplying that identity by $C$ gives \eqref{eq:radial-lck-periodic-ode}. Integration of $C^mC'$ gives the first identity in \eqref{eq:radial-lck-period-identities}, which forces $k<0$. The formulas for $l$ and $L$ follow from the Lee formula.
\end{proof}

\begin{lemma}\label{lem:radial-lck-periodic-order}
For fixed $F>0$, a positive two-periodic solution of \eqref{eq:radial-lck-periodic-ode} is unique for each $k<0$. If $k_2>k_1$, the corresponding solutions satisfy $C_2>C_1$ everywhere and $L(C_2)>L(C_1)$.
\end{lemma}
\begin{proof}
At a zero of $R=C_2-C_1$ with $k_2>k_1$, the equation gives $R'=(k_2-k_1)C_1^{1-m}>0$. Such zeros are isolated, and there are finitely many on the circle. A periodic differentiable function cannot have only strict upward crossings: returning around the circle would require a downward crossing. Hence $R$ has a constant strict sign. The identity $\int_0^2C_i\,\mathrm dr=-\int_0^2F\,\mathrm dr/k_i$ determines that sign as positive. The formula for $L(C)$ in \eqref{eq:radial-lck-period-identities} then gives its strict ordering.

For equal $k$, uniqueness for the smooth scalar equation $C'=kC^{1-m}+FC^{-m}$ on $C>0$ identifies solutions which meet. Solutions which never meet are strictly ordered, contradicting their equal integrals. This proves uniqueness for fixed $k$.
\end{proof}

\begin{theorem}\label{thm:radial-gauduchon-lck-volume}
For every positive smooth two-periodic $F$ and every $L_0<0$, there is exactly one metric in \eqref{eq:radial-lck-family} which is Gauduchon, has Lee period $L_0$, and satisfies
\[
 \omega^n=nF(r)\alpha^m\wedge\beta.
\]
\end{theorem}
\begin{proof}
Fix $k<0$, put $c=-k$, and let $f_0=\min F>0$, $f_1=\max F$. The constant endpoints of $I_k=[f_0/c,f_1/c]$ are respectively a subsolution and a supersolution of $C'=C^{-m}(F-cC)$. Scalar comparison keeps its forward flow in this compact positive interval. Smooth ODE existence and uniqueness therefore give a continuous time-two map $I_k\to I_k$. Its difference from the identity has opposite weak signs at the endpoints, so the intermediate value theorem gives a fixed point. If $f_0=f_1$, the constant $C=f_0/c$ is the solution. In either case the periodic equation makes all endpoint derivatives agree, giving a smooth positive periodic solution. It is unique by \cref{lem:radial-lck-periodic-order}, and satisfies
\begin{equation}\label{eq:radial-lck-ode-barriers}
 \frac{f_0}{-k}\leq C_k(r)\leq\frac{f_1}{-k}.
\end{equation}
The period $L(k)=-\int_0^2F/C_k^{m+1}\,\mathrm dr$ is strictly increasing. It is continuous: on a compact neighbourhood of any $k_0<0$, \eqref{eq:radial-lck-ode-barriers} gives uniform positive lower and finite upper bounds. The smooth equation then bounds every radial derivative. Any sequence $k\to k_0$ has smoothly convergent subsequences, and the equation and uniqueness identify every limit with $C_{k_0}$. Thus $C_k\to C_{k_0}$ uniformly, which gives continuity of the integral.

The same barriers imply $L(k)\to0$ as $k\uparrow0$ and $L(k)\to-\infty$ as $k\to-\infty$. Hence exactly one $k$ has period $L_0$. Set $D=F/C_k^m$. The identities of \cref{lem:radial-lck-gauduchon-equation} give the required metric, volume and Lee period. Every radial candidate gives a positive periodic solution of the same equation, so the period ordering also proves uniqueness throughout the stated family.
\end{proof}

\begin{proposition}\label{prop:radial-lck-positive-twisted-potential}
The metric of \cref{thm:radial-gauduchon-lck-volume} has a positive smooth global function $U$ with $\omega=D_\lambda U$ in its actual Lee gauge. Its class in $H^{1,1}_{\mathrm{BC},\lambda}(H_n,\mathbb R)$ is zero. For two such metrics with positive potentials $U_i$, $i=1,2$, and cohomologous Lee forms $\lambda_2=\lambda_1+\mathrm da$, every fixed choice of the additive constant in $a$ gives the class-difference potential $\mathrm e^{-a}U_2-U_1$. In particular, \eqref{eq:lck-coupled-equation} determines the metric among radial candidates, even when its smooth difference potential is not assumed radial.
\end{proposition}
\begin{proof}
Write $l=kC^{-m}<0$ and $f(r)=\int_0^rl(s)\,\mathrm ds$ on the real radial cover. Then $f(r+2)=f(r)+L_0$. The solution of $U'-lU=C$ with
\[
 \begin{aligned}
 U(0)&=\frac{\mathrm e^{L_0}}{1-\mathrm e^{L_0}}
       \int_0^2\mathrm e^{-f(s)}C(s)\,\mathrm ds,\\
 U(r)&=\mathrm e^{f(r)}
       \left(U(0)+\int_0^r\mathrm e^{-f(s)}C(s)\,\mathrm ds\right)
 \end{aligned}
\]
is positive on $[0,2]$ and has $U(2)=U(0)$. The smooth periodic equation matches all derivatives at the endpoints, so $U$ descends to a positive smooth function. Using the operator of \eqref{eq:lck-gauge-operators}, its actual twisted Hessian is
\[
 \begin{split}
 D_\lambda U
 &=(U'-lU)\alpha+
   \bigl[U''-2lU'+(l^2-l')U\bigr]\beta\\
 &=C\alpha+(C'-lC)\beta=\omega.
 \end{split}
\]
This proves the zero twisted class. The gauge identity
$\mathrm e^{-a}D_{\lambda_2}U_2=D_{\lambda_1}(\mathrm e^{-a}U_2)$ gives the asserted difference potential for the chosen class identification. For radial candidates of the coupled system, equal volume fixes the same $F$, and equal Lee class fixes the same period. Apply \cref{thm:radial-gauduchon-lck-volume}. This concerns the metric and imposes no uniqueness assertion on arbitrary representatives of its twisted potential.
\end{proof}

\begin{proposition}\label{prop:radial-lck-non-vaisman-volume}
If $F$ is nonconstant, the metric of \cref{thm:radial-gauduchon-lck-volume} is not Vaisman.
\end{proposition}
\begin{proof}
The radial complex line is orthogonal to $\ker\mathrm dr\cap\ker\dc r$. On a complex radial orbit, use a logarithmic coordinate $w=x+\mathrm i y$, so $r=2x+\text{constant}$ and $\beta=\pi^{-1}\mathrm dx\wedge\mathrm dy$. Its Riemannian metric is $D\pi^{-1}(\mathrm dx^2+\mathrm dy^2)$, whence
\[
 |\mathrm dr|_\omega^2=\frac{4\pi}{D},\qquad
 |\lambda|_\omega^2=4\pi k^2\frac{C^{-2m}}{D}.
\]
A parallel Lee form has constant norm. Since $k\ne0$, the Vaisman condition would imply $D=Q C^{-2m}$ for a constant $Q>0$. At both a maximum and a minimum of $C$, the Gauduchon identity would give $kC^{m+1}+Q=0$. This equation has just one positive root, so the maximum and minimum coincide. Thus $C$, $D$ and $F=C^mD$ would all be constant, a contradiction.
\end{proof}

\begin{example}\label{ex:radial-gauduchon-non-vaisman}
For
\[
 \begin{gathered}
 0<\varepsilon<\min\left\{\frac12,
                \frac{(3/2)^{1-m}}{4\pi}\right\},
 \\
 C=1+\varepsilon\cos(\pi r),\qquad
 D=C'+C^{1-m},
 \end{gathered}
\]
the metric \eqref{eq:radial-lck-family} is Gauduchon and non-Vaisman, with Lee form $\lambda=-C^{-m}\mathrm dr$ and density $F=C^mC'+C$. The probability volume $\mathrm dV=\frac12\alpha^m\wedge\beta$ gives $\omega^n=2nF\,\mathrm dV$ and $\int_{H_n}F\,\mathrm dV=1$.
\end{example}
\begin{proof}
Since $1/2<C<3/2$ and $m\geq1$, we have
$D\geq(3/2)^{1-m}-\pi\varepsilon>\frac34(3/2)^{1-m}>0$.
The Gauduchon identity holds with $k=-1$, giving the Lee form and the positive density $F=C^mD$. Its values $F(0)=1+\varepsilon$ and $F(1)=1-\varepsilon$ differ, so \cref{prop:radial-lck-non-vaisman-volume} applies. The integral of $C^mC'$ over a period is zero and the average of $C$ is one. Fibre integration then gives $\int_{H_n}F\,\mathrm dV=1$.
\end{proof}

The nonconstant radial volume in this example is not invariant under the Lee translation of the fixed Vaisman metric $\alpha+\beta$. The new metric's own radial Lee flux is different: with the factor $4\pi$ omitted, it is $lF/D=-1$, in agreement with its Gauduchon condition. The radial classification settles this family within \cref{q:lck-gauduchon-volume}; uniqueness among arbitrary LCK candidates remains open.
